\label{pp:fwx}
\section{FWX (FirmWare eXchange)}

% FWX\footnote{Název FWX vznikl náhodně při vývoji.} je specializované komunikační rozhraní mezi
% emulovaným programem a emulátorem (respektive firmwarem). Toto rozhraní je založeno na modelu

Aby mohl emulovaný systém (např. BIOS) přistupovat k hardwarovým prostředkům mimo standardní
architekturu NES -- jako je souborový systém na SD kartě nebo konfigurace zařízení -- bylo nutné
implementovat komunikační most, nazvaný \textbf{FWX}. Toto rozhraní umožňuje emulovanému softwaru
bezpečně předávat požadavky hostitelskému firmwaru emulátoru, a to prostřednictvím trojice paměťově
mapovaných registrů (\textbf{FWXCTRL}, \textbf{FWXSTAT} a \textbf{FWXDATA}).

\begin{figure}[htbp]
    \centering
    \includesvg[inkscapelatex=false]{images/typical_communication.svg}
    \caption{Zjednodušené schéma typického průběhu komunikace přes rozhraní FWX.}
    \label{fig:simplified_typical_communication}
\end{figure}

Rozhraní FWX je navrženo jako jednoduchý, bajtově orientovaný synchronní protokol založený na modelu
\textbf{příkaz-odpověď} (\textit{command-response}). Veškerou komunikaci iniciuje výhradně emulovaný
program (BIOS).

Rámcování zpráv je řízeno pomocí kontrolního bitu \emph{SS} (Start/Stop) v řídicím registru
\textbf{FWXCTRL}. Příkazy a jejich datové argumenty jsou přenášeny přes datový registr
\textbf{FWXDATA}. K zachycení zapisovaných dat emulátorem dochází při explicitní změně stavu
hodinového bitu \emph{CLK} (data latching). Vzhledem k tomu, že je protokol taktován softwarově
samotným emulovaným programem, nejsou zde kladeny žádné striktní požadavky na časování mezi
jednotlivými hodinovými událostmi.

Z pohledu emulovaného programu je exekuce příkazů blokující. Po odeslání požadavku je běh
emulace dočasně pozastaven, dokud hostitelský firmware příkaz nezpracuje. Dokončení operace je
následně signalizováno příznakem \emph{S0} ve stavovém registru \textbf{FWXSTAT}.

\subsection{Iniciační a terminační sekvence (START/STOP)}

Komunikační průběh je striktně ohraničen stavy START a STOP, které generuje emulovaný program.

\begin{itemize}
    \item \textbf{START podmínka} je vyvolána změnou bitu \emph{SS} z logické 0 na logickou 1. Po
        jejím vydání emulátor nastaví stavový bit \emph{S0} na 1 (indikace probíhající komunikace) a
        zároveň vynuluje příznaky \emph{S1}, \emph{S2} a chybové kódy \emph{ER0..ER3}.
    \item \textbf{STOP podmínka} nastává při změně bitu \emph{SS} z 1 na 0. Emulátor vynuluje
        stavové bity \emph{S0}, \emph{S1} a \emph{S2}, přičemž bity \emph{ER0..ER3} si zachovají
        chybový kód poslední provedené operace pro případnou kontrolu ze strany BIOSu.
\end{itemize}

\subsection{Přenos dat a příkazů}

\begin{figure}[htbp]
    \centering
    \includesvg[width=\textwidth,height=\textheight,inkscapelatex=false]{images/data_transfer_procedure.svg}
    \caption{Detailní stavový diagram procesu odesílání dat z emulovaného programu do emulátoru.}
    \label{fig:data_transfer_procedure}
\end{figure}

Samotný přenos informací směrem k emulátoru probíhá po iniciaci START podmínky. Emulovaný
program zapíše datový byte do registru \textbf{FWXDATA} a následně nastaví bit \emph{CLK} v registru
\textbf{FWXCTRL} na logickou 1. Jakmile emulátor data zpracuje, okamžitě bit \emph{CLK} automaticky
resetuje na 0, čímž potvrdí úspěšný příjem.

První byte odeslaný po START podmínce je vždy interpretován jako \textbf{ID příkazu} (Command ID).
Pokud daný příkaz vyžaduje dodatečné argumenty (např. název souboru nebo index), musí být tyto byty
odeslány bezprostředně po ID příkazu. Dokud emulátor neobdrží očekávaný počet datových bytů,
udržuje stavový bit \emph{S2} na hodnotě 1, čímž si žádá další data.

\subsection{Zpracování příkazu (blokující exekuce)}

Po úspěšném přijetí ID příkazu (a všech případných argumentů) zahájí hostitelský firmware jeho
exekuci. Během této fáze je emulace pozastavena. Tím je zaručena synchronizace --
emulovaný program se po provedení instrukce zápisu probudí přesně v okamžiku, kdy je operace na
straně emulátoru dokončena. Výsledek operace je reflektován prostřednictvím chybových kódů v
registru \textbf{FWXSTAT} (viz sekce \ref{subsec:last_error_code}).

\subsection{Příjem návratových dat}

\begin{figure}[htbp]
    \centering
    \includesvg[width=\textwidth,height=\textheight,inkscapelatex=false]{images/data_receive_procedure.svg}
    \caption{Detailní stavový diagram procesu příjmu návratových dat emulovaným programem.}
    \label{fig:data_receive_procedure}
\end{figure}

Některé příkazy (např. čtení obsahu adresáře) generují data, která musí být předána zpět emulovanému
programu. Přítomnost těchto návratových dat je signalizována nastavením bitu \emph{S1} na hodnotu 1.

Proces příjmu je inverzní k odesílání: emulovaný program si přečte hodnotu z registru
\textbf{FWXDATA} a pro vyžádání dalšího bajtu musí zapsat logickou 1 do bitu \emph{CLK}. Po obdržení
dalšího bytu, bit \emph{CLK} je automaticky nastaven na 0. Jakmile jsou všechna návratová data
úspěšně přečtena, emulátor nastaví bit \emph{S1} zpět na 0.

\subsection{Řetězení příkazů}

Komunikační protokol je navržen s ohledem na efektivitu; pokud je exekuce příkazu dokončena a na
lince je stále aktivní úroveň \emph{SS} (hodnota 1), může emulovaný program plynule navázat
odesláním dalšího ID příkazu bez nutnosti opětovně provádět STOP a START sekvenci.

\section{Mapování a přehled registrů}

\begin{table}[htbp]
    \centering
    \begin{tabular}{llcp{7cm}}
        \toprule
        \textbf{Adresa} & \textbf{Název} & \textbf{Směr} & \textbf{Účel} \\
        \midrule
        \texttt{\$4018} & \texttt{FWXSTAT} & R & Stavový registr a chybové kódy. \\
        \texttt{\$4019} & \texttt{FWXCTRL} & R/W & Řízení komunikace. \\
        \texttt{\$401A} & \texttt{FWXDATA} & R/W & Obousměrný datový registr. \\
        \bottomrule
    \end{tabular}
    \caption{Přehled komunikačních registrů rozhraní FWX.}
    \label{tab:fwx_regs}
\end{table}

Pro zajištění komunikace mezi emulovaným programem a emulátorem je rozhraní FWX mapováno přímo do
adresového prostoru CPU. V původním systému NES byl paměťový blok \texttt{\$4018-\$401F} vyhrazen
pro testovací režimy APU a CPU, avšak v produkčních verzích konzole byla tato funkčnost deaktivována
(tudíž není ani potřeba tuto funkčnost emulovat). Proto tento nevyužitý prostor je bezpečně
recyklován pro své komunikační registry, čímž předchází konfliktům se standardními hrami (je ovšem
nezbytné se vyhnout mapperům, které by tento prostor nestandardně využívaly).

Samotné rozhraní FWX se skládá ze tří 8bitových registrů, přičemž zbylý adresní prostor
(\texttt{\$401B-\$401F}) je rezervován pro budoucí rozšíření protokolu (např. plánované rozhraní pro
přímý přístup do paměti).

\subsection{Detailní popis registrů}

\label{sec:fwxstat}
\subsubsection{Stavový registr FWXSTAT (\texttt{\$4018})}

Tento registr slouží výhradně ke čtení ze strany emulovaného programu a poskytuje zpětnou vazbu o
stavu protokolu a případných chybách běhu. Výchozí hodnota po resetu je \texttt{0x00}.

\begin{figure}[H]
    \centering
    \includesvg{images/fwxstat_layout.svg}
    \caption{Rozvržení registru FWXSTAT.}
\end{figure}

\textbf{Stavové bity (S0--S2):}
\begin{itemize}
    \item \textbf{S0 (probíhá přenos)} je nastaven na logickou 1 po celou dobu trvání přenosu (mezi
        podmínkami START a STOP).
    \item \textbf{S1 (data připravena)} je nastaven na logickou 1 v případě, že emulátor po úspěšném
        zpracování příkazu vrací data, která je možné vyčíst z registru \texttt{FWXDATA}.
    \item \textbf{S2 (přenos argumentů)} indikuje, zda emulátor očekává další datové bajty pro
        aktuální příkaz. Dokud je \textbf{S2} rovno 0, probíhá sběr datových bajtů. Jakmile je
        obdržen potřebný počet bajtů a příkaz je proveden, je bit \textbf{S2} nastaven na 1. Příznak
        je automaticky vymazán při následném čtení registru nebo při vydání STOP podmínky.
\end{itemize}

Bit 7 je v současné verzi rozhraní rezervován.

\label{subsec:last_error_code}
\textbf{Chybové kódy (ER0--ER3):} spodní čtyři bity uchovávají chybový kód poslední provedené
operace. Emulovaný program by měl tyto bity kontrolovat po každém ukončeném příkazu. Seznam možných
chyb je uveden v tabulce \ref{tab:error_codes}.

\begin{table}[H]
    \centering
    \begin{tabular}{cll}
        \toprule
        \textbf{Hodnota} & \textbf{Mnemotechnika} & \textbf{Popis chyby} \\
        \midrule
        \texttt{0x0} & \texttt{ERR\_OK} & Operace proběhla úspěšně. \\
        \texttt{0x1} & \texttt{ERR\_BAD\_DATA} & Přijata neplatná struktura dat nebo neznámý příkaz. \\
        \texttt{0x2} & \texttt{ERR\_SD\_MISSING} & SD karta není vložena nebo selhala inicializace. \\
        \texttt{0x4} & \texttt{ERR\_SD\_IO\_ERROR} & Fyzická I/O chyba při práci se souborovým systémem. \\
        \texttt{0x5} & \texttt{ERR\_BAD\_IDX} & Index mimo rozsah v aktuálním pracovním adresáři. \\
        \texttt{0x6} & \texttt{ERR\_NON\_EXISTENT} & Požadovaný soubor nebo složka neexistuje. \\
        \texttt{0x7} & \texttt{ERR\_IS\_FILE} & Cílový objekt je soubor (očekáván adresář). \\
        \texttt{0x8} & \texttt{ERR\_IS\_DIR} & Cílový objekt je adresář (očekáván soubor). \\
        \bottomrule
    \end{tabular}
    \caption{Přehled chybových kódů vracených v registru FWXSTAT.}
    \label{tab:error_codes}
\end{table}

\label{sec:fwxctrl}
\subsubsection{Řídicí registr FWXCTRL (\texttt{\$4019})}

Slouží pro obousměrnou manipulaci s kontrolními signály linky. Zápisem do tohoto registru emulovaný
program fyzicky taktuje protokol. Výchozí hodnota je \texttt{0x00}.

\begin{figure}[htbp]
    \centering
    \includesvg{images/fwxctrl_layout.svg}
    \caption{Rozvržení registru FWXCTRL.}
\end{figure}

\begin{itemize}
    \item \textbf{SS (Start/Stop)} řídí životní cyklus komunikace. Zápis logické 1 otevírá relaci
        (START), zápis logické 0 relaci ukončuje (STOP).
    \item \textbf{CLK (Clock / Latch)} funguje jako softwarové hodiny protokolu.
        \begin{itemize}
            \item Při odesílání dat emulátoru (zápis do \texttt{FWXDATA}) se nastavením \textbf{CLK} na
                1 data \textit{uzamknou} (latch) a odešlou.
            \item Při čtení dat z emulátoru se nastavením \textbf{CLK} na 1 vyžádá další bajt odpovědi.
        \end{itemize}
    Změna bitu na 1 je ze strany emulátoru detekována ihned. Po zpracování (což zastaví cyklus
    procesoru) je bit \textbf{CLK} emulátorem automaticky vymazán zpět na 0, čímž je operace
    potvrzena a emulace může pokračovat.
\end{itemize}

Bity 2 až 7 jsou rezervovány.

\label{sec:fwxdata}
\subsubsection{Datový registr FWXDATA (\texttt{\$401A})}

Tento registr (čtení/zápis) slouží jako fyzické médium pro přesun samotných dat -- tedy ID příkazů,
jejich argumentů a návratových hodnot. Práce s tímto registrem je vždy podmíněna následnou
manipulací s bitem \textbf{CLK} v registru \texttt{FWXCTRL}. Výchozí hodnota je \texttt{0x00}.

\section{Sada příkazu}

Protokol FWX definuje sadu příkazů, pomocí kterých emulovaný program přistupuje k externím
prostředkům. Pro potřeby BIOSu byly implementovány především příkazy pro manipulaci se souborovým
systémem na SD kartě a pro řízení indikačních LED diod zařízení.

Veškerá data přenášená přes rozhraní FWX využívají následující datové typy:

\begin{itemize}
    \item \texttt{uint8\_t} -- standardní 8bitový bezznaménkový bajt.
    \item \texttt{cstring} -- textový řetězec znaků v kódování ASCII zakončený nulovým znakem
        (\texttt{0x00}), odpovídající standardu jazyka C.
\end{itemize}

\subsection{Přehled podporovaných příkazů}

Tabulka \ref{tab:fwx_cmds} shrnuje dostupné příkazy protokolu FWX. Očekávané argumenty
(\textit{Data}) musí být odeslány bezprostředně po ID příkazu. Návratové hodnoty (\textit{Návrat})
je nutné z emulátoru vyčíst před ukončením komunikace.

\begin{table}[htbp]
    \vspace{-2em}
    \centering
    \resizebox{\textwidth}{!}{
    \begin{tabular}{c l p{3.5cm} p{2cm} p{5.5cm}}
        \toprule
        \textbf{ID} & \textbf{Mnemotechnika} & \textbf{Argumenty (Data)} & \textbf{Návrat} & \textbf{Popis} \\
        \midrule
        \texttt{0x00} & \texttt{RESET} & \textit{žádné} & \textit{žádný} & Kompletní reinicializace
        emulátoru, vynulování RAM a znovunačtení BIOSu. \\
        \midrule
        \texttt{0x01} & \texttt{SD\_LOAD} & 1. \texttt{uint8\_t} (délka) \newline 2.
        \texttt{cstring} (cesta) & \textit{žádný} & Načte a spustí \texttt{.nes} ROM soubor. Pokud
        je cílem složka, změní do ní aktuální pracovní adresář. \\
        \midrule
        \texttt{0x02} & \texttt{SD\_LOAD\_IDX} & 1. \texttt{uint8\_t} (index) & \textit{žádný} &
        Analogické k \texttt{0x01}, ale cíl je specifikován indexem v aktuálním adresáři. \\
        \midrule
        \texttt{0x03} & \texttt{SD\_FILENAME\_IDX} & 1. \texttt{uint8\_t} (index) & \texttt{cstring}
        (název) & Vrátí název souboru na zadaném indexu. Neexistuje-li, vrátí prázdný řetězec. \\
        \midrule
        \texttt{0x04} & \texttt{SD\_IS\_DIR} & 1. \texttt{cstring} (cesta) & \texttt{uint8\_t}
        (bool) & Ověří, zda zadaná cesta ukazuje na složku (vrací 1) nebo soubor (vrací 0). \\
        \midrule
        \texttt{0x05} & \texttt{SD\_IS\_DIR\_IDX} & 1. \texttt{uint8\_t} (index) & \texttt{uint8\_t}
        (bool) & Analogické k \texttt{0x04}, ale cíl je specifikován indexem. \\
        \midrule
        \texttt{0x06} & \texttt{SD\_PWD} & \textit{žádné} & \texttt{cstring} (cesta) & Vrátí
        absolutní cestu k aktuálnímu pracovnímu adresáři emulátoru. \\
        \midrule
        \texttt{0x07} & \texttt{SD\_CWD} & 1. \texttt{cstring} (cesta) & \textit{žádný} & Změní
        aktuální pracovní adresář emulátoru na zadanou cestu. \\
        \midrule
        \texttt{0x08} & \texttt{SD\_FILE\_COUNT} & \textit{žádné} & \texttt{uint8\_t} (počet) &
        Vrátí celkový počet položek (souborů a složek) v aktuálním pracovním adresáři. \\
        \midrule
        \texttt{0x09} & \texttt{LED\_ERROR} & 1. \texttt{uint8\_t} (stav 1/0) & \textit{žádný} &
        Zapne (1) nebo vypne (0) červenou chybovou indikační LED na šasi zařízení. \\
        \midrule
        \texttt{0x0A} & \texttt{LED\_GENERAL} & 1. \texttt{uint8\_t} (stav 1/0) & \textit{žádný} &
        Zapne (1) nebo vypne (0) zelenou indikační LED (používáno např. při načítání). \\
        \bottomrule
    \end{tabular}
    }
    \caption{Přehled příkazů FWX rozhraní. Datové typy argumentů a návratových hodnot jsou uvedeny v
    závorkách.}
    \label{tab:fwx_cmds}
\end{table}

Poznámka k chybovým stavům: možné chyby (např. neexistující soubor \texttt{ERR\_NON\_EXISTENT} nebo
odpojená SD karta \texttt{ERR\_SD\_MISSING}) nejsou pro úsporu místa v tabulce \ref{tab:fwx_cmds}
explicitně vypsány. Emulovaný program musí po každém příkazu ověřit chybový kód v registru
\texttt{FWXSTAT}.
